\section{USAMO 1997}
        \begin{enumerate}
        	\item (\textit{USAMO 1997}) Let $p_1,p_2,p_3,...$ be the prime numbers listed in increasing order, and let $x_0$ be a real number between $0$ and $1$. For positive integer $k$, define


 $x_{k}=\begin{cases}0&\text{ if }x_{k-1}=0\\ \left\{\frac{p_{k}}{x_{k-1}}\right\}&\text{ if }x_{k-1}\ne0\end{cases}$


 where $\{x\}$ denotes the fractional part of $x$. (The fractional part of $x$ is given by $x-\lfloor{x}\rfloor$ where $\lfloor{x}\rfloor$ is the greatest integer less than or equal to $x$.) Find, with proof, all $x_0$ satisfying $0<x_0<1$ for which the sequence $x_0,x_1,x_2,...$ eventually becomes $0$.


 



\textbf{Solution}

All rational numbers between 0 and 1 inclusive will eventually yield some $x_k = 0$. To begin, note that by definition, all rational numbers can be written as a quotient of coprime integers. Let $x_0 = \frac{m}{n}$, where $m,n$ are coprime positive integers. Since $0<x_0<1$, $0<m<n$. Now 
 \[x_1 = \left\{\frac{p_1}{\frac{m}{n}}\right\}=\left\{\frac{np_1}{m}\right\}.\] From this, we can see that applying the iterative process will decrease the value of the denominator, since $m<n$. Moreover, the numerator is always smaller than the denominator, thanks to the fractional part operator. So we have a strictly decreasing denominator that bounds the numerator. Thus, the numerator will eventually become 0.

 On the other hand, if $x_0$ is irrational, then a simple induction will show that $x_k$ will always be irrational. Indeed, the base case has been established, and, if $x_k$ is irrational, then $\dfrac{p_{k+1}}{x_k}$ must be too, and likewise for its fractional portion, which differs from it by an integer. Hence $x_{k+1}$ is irrational, completing the proof.

 


	\item (\textit{USAMO 1997}) Let $ABC$ be a triangle, and draw isosceles triangles $BCD, CAE, ABF$ externally to $ABC$, with $BC, CA, AB$ as their respective bases. Prove that the lines through $A,B,C$ perpendicular to the lines $\overleftrightarrow{EF},\overleftrightarrow{FD},\overleftrightarrow{DE}$, respectively, are concurrent.


 



\textbf{Solution 1}

Let the perpendicular from A meet FE at A'. Define B' and C' similiarly. By \href{https://artofproblemsolving.com/wiki/index.php?title=Carnot\%27s\_Theorem}{Carnot's Theorem}, The three lines are \href{https://artofproblemsolving.com/wiki/index.php?title=Concurrent}{concurrent} if

 $FA'^2-EA'^2+EC'^2-DC'^2+DB'^2-FB'^2 = AF^2-AE^2+CE^2-CD^2+BD^2-BF^2 = 0$

 But this is clearly true, since D lies on the perpendicular bisector of BC, BD = DC.

 QED

 



\textbf{Solution 2}

We split this into two cases:

 Case 1: $D,E,F$ are non-collinear

 Observe that since $D,E,F$ lie on perpendicular bisectors, then we get that $DC=DB$, $EC=EA$, and $FA=FB$. This motivates us to construct a circle $\omega_1$ centered at $D$ with radius $DC$, and similarly construct $\omega_2$ and $\omega_3$ respectively for $E$ and $F$. 

 Now, clearly $\omega_1$ and $\omega_2$ intersect at $C$ and some other point. Now, we know that $DE$ is the line containing the two centers. So, the line perpendicular to $DE$ and passing through $C$ must be the radical axis of $\omega_1$ and $\omega_2$, which is exactly the line that the problem describes! We do this similarly for the others pairs of circles.

 Now by the radical lemma, the pairwise radical axes of $\omega_1,\omega_2,\omega_3$ are concurrent, as desired, and they intersect at the radical center.

 Case 2: $D,E,F$ are collinear

 Now, we are drawing perpendicular lines from $A$, $B$, and $C$ onto the single line $DEF$. Clearly, these lines are parallel and concur at the point of infinity.

 



\textbf{Solution 3}

Notice that $DEF$ and $ABC$ are orthologic, so the perpendiculars concur at the opposite center of orthology.

 


	\item (\textit{USAMO 1997}) Prove that for any integer $n$, there exists a unique polynomial $Q$ with coefficients in $\{0,1,...,9\}$ such that $Q(-2)=Q(-5)=n$.


 



\textbf{Solution}

The problem implies that $Q(x) = (x+2)(x+5) \cdot H_y(x) + n$. We define $H_y(x) = a_0 + a_1x + ... + a_nx^n$. Then, after expanding, the coefficient of $x^k$ is $10a_k + 7a_{k-1} + a_{k-2}$ where $2 \le k \le n$. The constant term is $10a_0 + n$. Since this has to be within the set ${0,1,...,9}$, $a_0$ must be unique according to $n$. The linear term is determined by $10a_1 + 7a_0$, which also must be within the set ${0,1,...,9}$. Since $a_0$ is determined uniquely by $n$, in order to keep the linear coefficient in the set ${0,1,...,9}$, $a_1$ is also determined uniquely. With $a_0$ and $a_1$ determined uniquely, we move on to the general coefficient. 

 Note that $10a_k + 7a_{k-1} + a_{k-2}$ is in the set ${0,1,...,9}$. We define $J_h(x)$ to be a function determining the remainder when $x$ is divided by $10$. Thus, for all $x$, $0 \le J_h(x) \le 9$. Note also that $J_h(10a_k + 7a_{k-1} + a_{k-2}) = 10a_k + 7a_{k-1} + a_{k-2}$, since $10a_k + 7a_{k-1} + a_{k-2}$ is in the range of $J_h(x)$. Therefore, $J_h(7a_{k-1} + a_{k-2}) = 10a_k + 7a_{k-1} + a_{k-2}$, which results in $a_k = \lfloor \frac{7a_{k-1} + a_{k-2}}{10} \rfloor$ for all $2 \le k \le n$. Thus, $a_k$ is determined uniquely, since all coefficients $a_i$ where $0 \le i < k$ are uniquely determined, which in turn determines a unique polynomial $Q(x)$.

 


	\item (\textit{USAMO 1997}) To \textit{clip} a convex $n$-gon means to choose a pair of consecutive sides $AB, BC$ and to replace them by three segments $AM, MN,$ and $NC,$ where $M$ is the midpoint of $AB$ and $N$ is the midpoint of $BC$. In other words, one cuts off the triangle $MBN$ to obtain a convex $(n+1)$-gon. A regular hexagon $P_6$ of area $1$ is clipped to obtain a heptagon $P_7$. Then $P_7$ is clipped (in one of the seven possible ways) to obtain an octagon $P_8$, and so on. Prove that no matter how the clippings are done, the area of $P_n$ is greater than $\frac{1}{3}$, for all $n\ge6$.


 



\textbf{Solution}


                    \begin{figure}[H]
                        \centering
                        \includegraphics[width=0.5\textwidth]{Solution/USAMO_I_1997_4_Solution_Figure_1.png}
                    \end{figure}
                    

 $\textbf{Claim:}$  It is impossible to choose two non-adjacent sides and clip a whole part of it off.

 $\textbf{Proof:}$  If you clip adjacent sides, you can cut off at most up to the blue lines; Clipping more is impossible due to the degrees getting larger and larger and more and more circular.

 Thus, after infinitely many clips, the hexagon bounded by the blue lines is left, so after finitely many clips, the area left is more than that hexagon.

 The side length of that hexagon is $\frac{\sqrt{3}}{3}$ of the large one, because of 30-30-120 triangles. Thus, the area is $\frac13$ of the larger one, so we are done.

 


	\item (\textit{USAMO 1997}) Prove that, for all positive real numbers $a, b, c,$


 $\frac{1}{a^3+b^3+abc}+\frac{1}{b^3+c^3+abc}+\frac{1}{a^3+c^3+abc} \le \frac{1}{abc}$.


 



\textbf{Solution 1}

Because the inequality is homogenous (i.e. $(a, b, c)$ can be replaced with $(ka, kb, kc)$ without changing the inequality other than by a factor of $k^n$ for some $n$), without loss of generality, let $abc = 1$.

 Lemma:
 \[\frac{1}{a^3 + b^3 + 1} \le \frac{c}{a + b + c}.\]Proof: Rearranging gives $(a^3 + b^3) c + c \ge a + b + c$, which is a simple consequence of $a^3 + b^3 = (a + b)(a^2 - ab + b^2)$ and
 \[(a^2 - ab + b^2)c \ge (2ab - ab)c = abc = 1.\]

 Thus, by $abc = 1$:
 \[\frac{1}{a^3 + b^3 + abc} + \frac{1}{b^3 + c^3 + abc} + \frac{1}{c^3 + a^3 + abc}\]
 \[\le \frac{c}{a + b + c} + \frac{a}{a + b + c} + \frac{b}{a + b + c} = 1 = \frac{1}{abc}.\]

 



\textbf{Solution 2}

Rearranging the AM-HM inequality, we get $\frac{1}{x}+\frac{1}{y}+\frac{1}{z} \le \frac{9}{x+y+z}$. Letting $x=a^{3}+b^{3}+abc$, $y=b^{3}+c^{3}+abc$, and $z=c^{3}+a^{3}+abc$, we get 
 \[\frac{1}{a^{3}+b^{3}+abc}+\frac{1}{b^{3}+c^{3}+abc}+\frac{1}{c^{3}+a^{3}+abc} \le \frac{9}{2a^{3}+2b^{3}+2c^{3}+3abc}.\] By AM-GM on $a^{3}$, $b^{3}$, and $c^{3}$, we have 
 \[a^{3}+b^{3}+c^{3} \ge 3abc \Rightarrow 2a^{3}+2b^{3}+2c^{3}+3abc \ge 9abc \Rightarrow \frac{9}{2a^{3}+2b^{3}+2c^{3}+3abc} \le \frac{1}{abc}.\] So, $\frac{1}{a^{3}+b^{3}+abc}+\frac{1}{b^{3}+c^{3}+abc}+\frac{1}{c^{3}+a^{3}+abc} \le \frac{1}{abc}$.-Tigerzhang

 $\textbf{WARNING:}$

 This solution doesn’t work because $(x+y+z)\left(\frac{1}{x}+\frac{1}{y}+\frac{1}{z}\right)\geq 9$, so $\frac{1}{x}+\frac{1}{y}+\frac{1}{z}\geq \frac{9}{x+y+z}$

 



\textbf{Solution 3}

If we multiply each side by $abc$, we get that we must just prove that 
 \[\frac{abc}{a^3+b^3+abc} + \frac{abc}{b^3+c^3+abc} + \frac{abc}{c^3+a^3+abc} \leq 1\] If we divide our LHS equation, we get that 
 \[\frac{1}{\frac{a^3+b^3+abc}{abc}} +\frac{1}{\frac{b^3+c^3+abc}{abc}} + \frac{1}{\frac{c^3+a^3+abc}{abc}}\] 
 \[= \frac{1}{\frac{a^2}{bc} + \frac{b^2}{ac} + 1} +\frac{1}{\frac{b^2}{ac} + \frac{c^2}{ab} + 1}  + \frac{1}{\frac{c^2}{ab} + \frac{a^2}{bc} + 1}\] Make the astute observation that by Titu's Lemma, 
 \[\frac{a^2}{bc} + \frac{b^2}{ac} \geq \frac{\left(a+b\right)^2}{bc+ ac}\] 
 \[\sum_{cyc}\frac{\left(a+b\right)^2}{bc+ ac} \leq \sum_{cyc}\frac{a^2}{bc} + \frac{b^2}{ac}\] Therefore: 
 \[\sum_{cyc} \frac{1}{\frac{\left(a+b\right)^2}{c\left(a+b\right)  }+1} \geq \sum_{cyc} \frac{1}{\frac{a^2}{bc} + \frac{b^2}{ac} + 1}\] If we expand it out, we get that 
 \[\sum_{cyc} \frac{1}{\frac{\left(a+b\right)^2}{c\left(a+b\right)  }+1}  = \frac{a+b+c}{a+b+c} = 1\] Since our original equation is less than this, we get that 
 \[\sum_{cyc} \frac{1}{\frac{a^2}{bc} + \frac{b^2}{ac} + 1} \leq 1\] 
 \[\sum_{cyc} \frac{abc}{a^3+b^3+abc}\leq 1\] 
 \[\sum_{cyc} \frac{1}{a^3+b^3+abc} \leq \frac{1}{abc}\]

 -KEVIN\_LIU

 


	\item (\textit{USAMO 1997}) Suppose the sequence of nonnegative integers $a_1,a_2,...,a_{1997}$ satisfies 


 $a_i+a_j \le a_{i+j} \le a_i+a_j+1$


 for all $i, j \ge 1$ with $i+j \le 1997$. Show that there exists a real number $x$ such that $a_n=\lfloor{nx}\rfloor$ (the greatest integer $\le x$) for all $1 \le n \le 1997$.


 



\textbf{Solution}

Given nonnegative integers $a_1,a_2,\dots, a_{1997}$ satisfying the given inequalities, let $I_n$ be the set of all $x$ such that $a_n=\lfloor nx\rfloor$. Therefore,
 \[I_n=\{x\,:\,a_n\le nx<a_n+1\}.\]We can rewrite this as
 \[I_n=\left[\frac{a_n}{n},\frac{a_n+1}{n}\right).\tag{1}\]So we know that each $I_n$ is an interval. We start by proving that if $n$ and $k$ are positive integers, then $I_n\cap I_{k}\ne \emptyset$. To prove this, we need the following lemma.

 \textbf{Lemma 1:} If $n$ and $k$ are positive integers, then
 \[na_k< k(a_n+1)\tag{2}\]

 \textbf{Proof:} We prove this by induction. If $k=1$, then we wish to show that $na_1< a_n+1$. By the given lower bound, we know that
 \[a_n\ge a_{n-1}+a_1\ge (a_{n-2}+a_1)+a_1\ge \cdots \ge na_1.\]Therefore, $na_1\le a_n<a_n+1$, so the hypothesis is true for $k=1$. If $n=1$, then we wish to prove that $a_k < k(a_1+1)$. By the given upper bound, we know that
 \[a_k\le a_{k-1}+(a_1+1)\le (a_{k-2}+(a_1+1))+(a_1+1)\le\cdots\le ka_1+(k-1).\]Therefore, $a_k<ka_1+k=k(a_1+1)$, so the hypothesis is true for $n=1$.

 Now suppose that (2) holds for all $k,n<j$. If $k=n=j$, then (2) is equivalent to $0<j$, which is obviously true. If $n=j>k$, then by the division algorithm, we can write $n=qk+r$ for nonnegative integers $q$ and $r$ with $0\le r<k$. Thus by applying the lower bound repeatedly (with one of the indices equal to $1$), we find
 \[k(a_n+1)\ge k(qa_k+a_r+1)=kqa_k+k(a_r+1).\tag{3}\]But then as $k,r<j$, we can apply the inductive hypothesis with $n=r$ and $k=k$ to find $k(a_r+1)> ra_k$. Substituting this into (3), we find
 \[k(a_n+1)> kqa_k+ra_k=(kq+r)a_k=na_k.\]This proves the hypothesis for $n=j>k$.

 Suppose that $k=j>n$. Then by the division algorithm, we can write $k=qn+r$ for nonnegative integers $q$ and $r$ with $0\le r<n$. Then by applying the upper bound repeatedly (with one of the indices equal to $1$), we find
 \[na_k\le n(qa_n+a_r+q)=nqa_n+na_r+nq.\tag{4}\]But as $n,r<j$, we can apply the inductive hypothesis with $n=n$ and $k=r$, finding that $na_r< r(a_n+1)$. Subsituting this into (4), we find
 \[na_k< nqa_n+r(a_n+1)+nq=(nq+r)(a_n+1)=k(a_n+1).\]This proves the hypothesis for $k=j>n$.

 Therefore, if the hypothesis is true for all $k,n<j$, then it must be true for all $k,n\le j$. Hence by induction, it must be true for all positive integers $k$ and $n$. $\hfill\ensuremath{\square}$

 Now suppose for the sake of contradiction that $I_n$ and $I_k$ are disjoint intervals. Without loss of generality, we may assume that $I_n$ precedes $I_k$ on the number line. Hence the upper bound of $I_n$ is less than or equal to the minimum value of $I_k$, or rather
 \[\frac{a_{n}+1}{n}\le \frac{a_k}{k}.\]This simplifies to $na_k\ge k(a_n+1)$. But this contradicts the statement of Lemma 1. Therefore, $I_k\cap I_n\ne \emptyset$.

 Now we claim that $I_1\cap I_2\cap \cdots \cap I_{1997}\ne \emptyset$. We prove this by induction. By the above, we know that $I_1\cap I_2\ne \emptyset$. Now suppose that $I_1\cap I_2\cap\cdots\cap I_k\ne \emptyset$. The intersection of two overlapping intervals of the form $[a,b)$ and $[c,d)$ is an interval of the form $[e,f)$, where $e=\max\{a,c\}$ and $f=\min\{b,d\}$. Therefore, by induction, we know that if the intersection of $k$ overlapping intervals is nonempty, then it must also be an interval, say
 \[I_1\cap I_2\cap\cdots\cap I_k=[x_k,y_k).\]If $I_{k+1}$ does not intersect $[x_k,y_k)$, then as an interval, it must appear either completely before $x_k$ or completely after $y_k$.  If $I_{k+1}$ appears completely before $x_k$, then it has a nonempty intersection with each of $I_1, I_2, \dots, I_k$. But we also know that $x_k$ is a lower bound of one of the intervals, hence $I_{k+1}$ cannot intersect that interval, a contradiction. A similar contradiction arises if $I_{k+1}$ appears completely after $y_k$. Therefore,
 \[I_1\cap I_2\cap\cdots\cap I_{k+1}\ne \emptyset.\]Thus by induction, $I_1\cap I_2\cap\cdots\cap I_{1997}\ne \emptyset$. So let $x\in I_1\cap I_2\cap\cdots\cap I_{1997}$. Then by the definition of $I_n$, we know that $a_n=\lfloor nx\rfloor$ for all $1\le n\le 1997$, and we are done.

 


\end{enumerate}